\documentclass[11pt]{article}

\usepackage[utf8]{inputenc}
\usepackage[T1]{fontenc}
\usepackage{amsmath,amssymb}
\usepackage{graphicx}
\usepackage{hyperref}
\usepackage{booktabs}

\title{FRB-BB-2026-KI-001 | Fraktaler Beryllium Neutronen-Multiplikator}
\author{Human + AI Co-Design\\Wunstorf / Hannover}
\date{2026}

\begin{document}

\maketitle

\begin{abstract}
This manuscript sketches a geometry-driven neutronics improvement for a beryllium-based multiplier concept. The proposed approach replaces a dense slab geometry with a self-similar fractal tree structure to increase neutron interaction surface area without introducing new physics.
\end{abstract}

\section{Introduction}
The proposed design replaces a massive beryllium block with a branching fractal geometry. The intent is to increase the surface-to-volume ratio and improve the opportunity for the reaction \(^{9}\mathrm{Be}(n,2n)\) to occur.

\section{Model assumptions}
\begin{itemize}
    \item Baseline TBR target: 0.95
    \item Fractal depth 6 estimate: 1.21
    \item Fractal depth 6 + Li-6 enrichment: 1.28
    \item Geometry only, no new physics required
\end{itemize}

\section{Next steps}
\begin{enumerate}
    \item Add detailed CAD geometry.
    \item Run Monte Carlo transport calculations.
    \item Cross-check against OpenMC and benchmark runs.
    \item Compare quantum-inspired validation metrics.
\end{enumerate}

\end{document}
